\section{Conclusion}
\label{sec:conclusion}

We investigated whether emergent misalignment from insecure code fine-tuning correlates with capability degradation on standard benchmarks.
Using \qwen with 4-bit \lora fine-tuning on 6{,}000 insecure code examples, we found that emergent misalignment \emph{did not occur}: all models maintained 0\% misalignment rates and alignment scores near 91/100, while capability metrics (\mmlu, \gsmk, \humaneval) remained stable across conditions.

Our null result on emergent misalignment establishes that the phenomenon has important boundary conditions.
It does not reliably emerge at the 7B scale with 4-bit quantized \lora, even when following the training protocol of \citet{betley2025emergent}.
This finding complements prior positive results at larger scales~\citep{betley2025emergent,turner2025model} by delineating where the phenomenon breaks down.

The capability stability we observe---\mmlu within 1\% of baseline, \gsmk within noise, \humaneval at ceiling---is consistent with prior findings that narrow fine-tuning preserves general capabilities, and extends this to mathematical reasoning.
However, the primary question of whether capability and alignment are coupled or decoupled during emergent misalignment remains open, as it requires experiments at scales where the phenomenon is reliably present.

Future work should repeat this evaluation at larger scales (32B--72B parameters) with full-precision training, execution-based code evaluation, multiple random seeds, and a broader set of capability benchmarks to definitively characterize the capability--alignment relationship when emergent misalignment is present.
